%!TEX root = ../main.tex

\titlespacing*{\section}{0pt}{2.0ex}{2.0ex}
\titlespacing*{\subsection}{0pt}{2.0ex}{2.0ex}

{\large\textbf{Aaron Mishkin}}
\hfill
\begin{tabular}{ll}
    %\normalsize \today\\[-.5em]
    \normalsize \href{mailto:aaronpmishkin@gmail.com}{\texttt{aaronpmishkin@gmail.com}} \\[-.25em]
    \normalsize \href{https://cs.stanford.edu/~amishkin/}{\texttt{https://cs.stanford.edu/{\textasciitilde}amishkin/}}
\end{tabular}

%--------------------SECTIONS-----------------------------------
\newenvironment{fourtable}{%
    \begin{tabular}{p{.15\textwidth}p{.445\textwidth}p{.145\textwidth}p{.15\textwidth}}%
}{%
    \end{tabular}%
}

\newenvironment{threetable}{%
    \begin{tabular}{p{.62\textwidth}p{.145\textwidth}p{.155\textwidth}}%
}{%
    \end{tabular}%
}
\newenvironment{threetableother}{%
    \begin{tabular}{p{.205\textwidth}p{.56\textwidth}p{.14\textwidth}}%
}{%
    \end{tabular}%
}
\newenvironment{twotable}{%
    \begin{tabular}{p{.205\textwidth}p{.72\textwidth}}%
}{%
    \end{tabular}%
}

\vspace{2ex}
\subsection{Employment}

\begin{description}[leftmargin=\mleftsize,style=nextline,font=\normalfont\space]
    \mitem{2026--} 
    Postdoctoral Researcher \hfill (starting 2026--10--01)\\
    École Polytechnique Fédérale de Lausanne (EPFL), Switzerland\\
    Advisor: Martin Jaggi
\end{description}

\subsection{Education}

\begin{description}[leftmargin=\mleftsize,style=nextline,font=\normalfont\space]
    \mitem{2026}
    PhD. in Computer Science \hfill Oral Defense 2025--11--18\\
    Stanford University, USA \hfill Degree conferred 2026--04--02\\
    Thesis: Convex Analysis of Non-Convex Neural Networks\\
    Advisor: Mert Pilanci

    \mitem{2020}
    MSc. in Computer Science \hfill Degree conferred 2020--11--01\\
    University of British Columbia (UBC), Canada	\\
    Thesis: Interpolation, Growth Conditions, and Stochastic Gradient Descent \\
    Advisor: Mark Schmidt

    \mitem{2018} 
    BSc. in Computer Science \hfill Degree conferred 2018--05--01\\
    University of British Columbia (UBC), Canada \\
    Advisor: David Poole
\end{description}

\subsubsection{Research Experience}
\begin{description}[leftmargin=\mleftsize,style=nextline,font=\normalfont\space]
    \item[2025]
          \textbf{Research Intern}, The Voleon Group, USA  \hfill (2025--05--13 to 2025--08--01)
          \\
          Advisors: Sahand Negahban and Deepak Rajan
    \item[2024]
          \textbf{Visiting PhD Student}, Inria Sierra Team, France \hfill (2024--07--01 to 2024--12--31) 
          \\
          Advisor: Francis Bach 
    \item[2023]
          \textbf{Predoctoral Researcher}, Flatiron Institute, USA \hfill (2023--06--01 to 2023--08--11) 
          \\
          Advisors: Robert Gower and Alberto Bietti 
    \item[2019]
          \textbf{Applied Science Intern}, Amazon Development Centre Berlin, Germany \hfill (2019--05--01 to 2019--08--31) 
          \\
          Advisors: Cédric Archambeau and Matthias Seeger
    \item[2018]
          \textbf{Research Intern}, RIKEN AIP Centre, Japan \hfill (2018--01--01 to 2019--06--31) 
          \\
          Advisors: Emtiyaz Khan 
\end{description}

\subsubsection{Other Experience}
\begin{description}[leftmargin=\mleftsize,style=nextline,font=\normalfont\space]
    \item[2017]
          \textbf{NSERC Undergraduate Research Assistant}, UBC, Canada  \hfill (2017--05--01 to 2017--08--31)
          \\
          Advisors: David Poole and Giuseppe Carenini
    \item[2016]
          \textbf{NSERC Undergraduate Research Assistant}, UBC, Canada  \hfill (2016--05--01 to 2016--08--31)
          \\
          Advisors: David Poole and Giuseppe Carenini
    \item[2015]
          \textbf{Co-op Student}, MacDonald, Detwiler, and Associates, Canada \hfill (2015--05--01 to 2015--12--31) 
\end{description}

\vspace{2ex}

\subsection{Publications
    \texorpdfstring{\hfill {\normalfont\small * denotes equal contribution.}}{}
}

\subsubsection{In preparation/submission}
\begin{description}[leftmargin=\mleftsize,style=nextline,font=\normalfont\space]
    \mitem{2026}
    \textit{Analyzing and Improving Greedy 2-Coordinate Updates for
        Equality-Constrained Optimization via Steepest Descent in the 1-Norm.}\\
    A. V. Ramesh$^{*}$, \textbf{A. Mishkin}$^{*}$, M. Schmidt, Y. Zhou, J. W.
    Lavington, and J. She\\
    In preparation.
    %\href{https://arxiv.org/abs/2402.19449}{arxiv:2402.19449}
\end{description}

\subsubsection{Refereed Papers}
\begin{description}[leftmargin=\mleftsize,style=nextline,font=\normalfont\space]
    \mitem{2026\\[.25em] TMLR}
    \textit{Glocal Smoothness: Line search and adaptive sizes can help in theory too!}\\
    C. Fox, \textbf{A. Mishkin}, S. Vaswani, M. Schmidt\\
    Transactions on Machine Learning Research (TMLR) 2026

    \mitem{2026\\[.25em] SIMODS}
    \textit{A Library of Mirrors: Deep Neural Nets in Low Dimensions are Convex
        Lasso Models with Reflection Features.}\\
    E. Zeger, Y. Wang, \textbf{A. Mishkin}, T. Ergen, E. Candès, M.  Pilanci\\
    SIAM Journal on Mathematics of Data Science (SIMODS) 2026

    \mitem{2025\\[.25em] AISTATS}
    \textit{Level Set Teleportation: An Optimization Perspective}\\
    {\bf A. Mishkin}, A. Bietti, R. M. Gower\\
    International Conference on Artificial Intelligence and Statistics (AISTAT) 2025
    \hfill 

    \mitem{2025\\[.25em] ICLR}
    \textit{Exploring the Loss Landscape of Regularized Neural Networks via Convex Duality}
    \hfill {\bf Oral}
    \\
    S. Kim, {\bf A. Mishkin}, M. Pilanci\\
    International Conference on Learning Representations (ICLR) 2025

    \mitem{2024\\[.25em] NeurIPS}
    \textit{Directional Smoothness and Gradient Methods: Convergence and Adaptivity}\\
    {\bf A. Mishkin}$^*$, A. Khaled$^*$, Y. Wang, A. Defazio, R. M. Gower\\
    Advances in Neural Information Processing Systems (NeurIPS) 2024

    \mitem{2023\\[.25em] ICML}
    \textit{Optimal Sets and Solution Paths of ReLU Networks}\\
    {\bf A. Mishkin}, M. Pilanci\\
    International Conference on Machine Learning (ICML) 2023

    \mitem{2022\\[.25em] ICML}
    \textit{Fast Convex Optimization for Two-Layer ReLU Networks}\\
    {\bf A. Mishkin}, A. Sahiner, M. Pilanci\\
    International Conference on Machine Learning (ICML) 2022

    \mitem{2019\\[.25em] NeurIPS}
    \textit{Painless Stochastic Gradient: Interpolation, Line-Search, and Convergence Rates}\\
    S. Vaswani, {\bf A. Mishkin}, I. H. Laradji, M. Schmidt, G. Gidel, S. Lacoste-Julien\\
    Advances in Neural Information Processing Systems (NeurIPS) 2019

    \mitem{2018\\[.25em] NeurIPS}
    \textit{SLANG: Fast Structured Covariance Approximations for Bayesian Deep Learning with Natural Gradient}\\
    {\bf A. Mishkin*}, F. Kunstner*, D. Nielsen, M. Schmidt, M.E. Khan\\
    Advances in Neural Information Processing Systems (NeurIPS) 2018
\end{description}

\vspace{1ex}

\subsubsection{Book Chapters}

\begin{description}[leftmargin=\mleftsize,style=nextline,font=\normalfont\space]
    \mitem{2022\\[.25em]} 
    %Contribution to book chapter\\
    \textit{Probabilistic Machine Learning: An Introduction}, K.P. Murphy, MIT Press, March 2022.\\
    Contributors: F. Kunstner, S.Y. Meng, \textbf{A. Mishkin}, S. Vaswani, M. Schmidt.\\
    Relevant Section: Chapter 8, \textit{Optimization}.
\end{description}

\vspace{1ex}

\subsubsection{Workshop Papers}

\begin{description}[leftmargin=\mleftsize,style=nextline,font=\normalfont\space]
    \mitem{2023\\[.25em] OptML}
    \textit{A Novel Analysis of Gradient Descent under Directional Smoothness}\\
    {\bf A. Mishkin}$^*$, A. Khaled$^*$, A. Defazio, R. M. Gower\\
    {NeurIPS Workshop on Optimization for Machine Learning} 2023

    \mitem{2023\\[.25em] OptML}
    \textit{Level Set Teleportation: the Good, the Bad, and the Ugly}\\
    \textbf{A Mishkin}, A. Bietti, R. M. Gower\\
    {NeurIPS Workshop on Optimization for Machine Learning} 2023
    % (\href{https://openreview.net/forum?id=7Ie3aoqJg1}{openreview:7Ie3aoqJg1})

    \mitem{2022\\[.25em] OptML}
    {\it The Solution Path of the Group Lasso}\\
    \textbf{A. Mishkin}, M. Pilanci\\
    {NeurIPS Workshop on Optimization for Machine Learning} 2022

    \mitem{2022\\[.25em] OptML}
    {\it Fast Convergence of Greedy 2-Coordinate Update for Optimizing with an Equality Constraint}\\
    A. V. Ramesh, {\bf A. Mishkin}, Mark Schmidt\\
    {NeurIPS Workshop on Optimization for Machine Learning} 2022
    % (\href{https://arxiv.org/abs/2304.13960}{arxiv:2304.13960})

    \mitem{2020\\[.25em] OptML}
    {\it To Each Optimizer a Norm, to Each Norm its Generalization}\\
    S. Vaswani, R. Babanezhad, J. Gallego, {\bf A. Mishkin}, S. Lacoste-Julien, N. Le Roux\\
    {NeurIPS Workshop on Optimization for Machine Learning} 2020
\end{description}

\vspace{2ex}
\subsection{Highlighted Research}

\begin{description}[leftmargin=\mleftsize,style=nextline,font=\normalfont\space]

    \mitem{1.}

    \textit{Optimal Sets and Solution Paths of ReLU Networks.}\\
    {\bf A.  Mishkin}, M. Pilanci, ICML 2023.

    This paper originated from a short proof showing that the solution path (the
    mapping from regularization parameter to optimal model) can be
    discontinuous for very narrow ReLU networks.
    This result is concerning because standard machine learning
    practice is to tune the regularization with grid-search, but 
    grid-search cannot have any theoretical guarantee when the solution path 
    is discontinuous.
    Trying to understand this behavior led us to study the structure of optimal
    two-layer ReLU using convex reformulations. 
    Using this approach, we show how to exactly characterize the optimal set
    for two-layer ReLU networks and, as
    a consequence, we also derive an optimal neuron pruning algorithm.
    This pruning method relies on a novel representation of neuron-sparse
    networks as extreme points of a ``optimal polyhedron''. 
    A follow-up project with Sungyoon Kim received an oral at ICLR for
    leveraging our work to understand mode connectivity.

    \mitem{2.}
    \textit{Directional Smoothness and Gradient Methods: Convergence and
        Adaptivity.} \\
    {\bf A. Mishkin}$^*$, A. Khaled$^*$, Y. Wang, A. Defazio, R. M. Gower, NeurIPS 2024

    The goal of this work is to bring the analysis of gradient methods down
    to the level of the objective function being optimized.
    To achieve this, we describe the convergence of gradient descent using
    \emph{directional smoothness functions}, which replace global Lipschitz
    continuity of the gradient with point-wise smoothness between the specific iterates
    generated by the optimizer.
    This allows us to study objectives which do not have Lipschitz gradients
    (ReLU networks, for example) and leads to rigorous convergence bounds which
    are an order of magnitude tighter than standard rates.
    These rates are also path-dependent, meaning that they capture the distinct
    behavior of each exact execution of gradient descent.
    Unlike other path-dependent approaches, our bounds can be evaluated
    in hindsight and used to compare algorithms.

    \mitem{3.}
    \textit{Painless Stochastic Gradient: Interpolation, Line-Search, and
        Convergence Rates} \\
    S. Vaswani, {\bf A. Mishkin}, I. H. Laradji, M. Schmidt, G. Gidel, S. Lacoste-Julien,
    NeurIPS 2019

    In this project, we circumvent hyper-parameter search
    using the surprising observation that powerful machine learning models often
    exactly fit, or \emph{interpolate}, their training data.
    In such cases, stochastic gradient descent shows a remarkable ``automatic
    variance reduction'' property and can converge with a fixed step-size.
    We show that SGD can also converge using a stochastic
    line-search (SLS), which looks along a random gradient for a
    step-size which leads to descent on average.
    We prove that SLS obtains fast rates of convergence under interpolation
    that match deterministic GD up to a constant factor.
    This analysis was among the first to use interpolation conditions in
    optimization and has been highly influential.

\end{description}

\vspace{2ex}
\subsection{Awards}

\subsubsection{Grants \& Fellowships}
\begin{description}[leftmargin=\mleftsize,style=nextline,font=\normalfont\space]
    % https://aistats.org/aistats2021/awards.html
    \item[2026]
          Canada Postgraduate Research Award Program (PDRA)
          \hfill (Declined) \$140,000 CAD\\
          Natural Sciences and Engineering Research Council (NSERC) of Canada
    \item[2026]
          \textit{Postdoctoral Research Associateship}
          \hfill (Declined) \pounds83,272 GBP\\
          Mathematical Institute, University of Oxford
    \item[2024] Visiting Student Researcher Fellowship, \hfill \$7 200 USD\\
          France-Stanford Center For Interdisciplinary Studies
    \item[2020] Graduate Research Fellowship Program (GRFP) \hfill \$138,000 USD\\
          National Sciences Foundation (NSF), USA
    \item[2020] Postgraduate Scholarships-Doctoral Program \hfill \$66,000 CAD\\
          Natural Sciences and Engineering Research Council (NSERC) of Canada
    \item[2020] Canada Postgraduate Scholarships-Doctoral Program \hfill (Declined) \$105,000 CAD\\
          Natural Sciences and Engineering Research Council (NSERC) of Canada
    \item[2019] Huawei Graduate Scholarship \hfill \$5,000 CAD\\
          Department of Computer Science, University of British Columba
    \item[2018] Computer Science Merit Scholarship \hfill \$20,000 CAD\\
          Department of Computer Science, University of British Columba
    \item[2018] Alexander Graham Bell Canada Graduate Scholarships-Master's \hfill \$17,500 CAD\\
          Natural Sciences and Engineering Research Council of Canada (NSERC)
\end{description}

\subsubsection{Scholarships \& Prizes}\nobreak
\begin{description}[leftmargin=\mleftsize,style=nextline,font=\normalfont\space]
    \item[2025] Best Poster\\
          ICCOPT 2025
    \item[2018] Academic Award of Excellence (Honors)\\
          Department of Computer Science, University of British Columba
    \item[2018] Markus Meister Memorial Prize\\
          Department of Computer Science, University of British Columba
    \item[2018] Travel Award\\
          NeurIPS Foundation
    \item[2017] D. F. MacKenzie Scholarship\\
          University of British Columba
    \item[2017] Undergraduate Student Research Award\\
          Natural Sciences and Engineering Research Council of Canada (NSERC)
    \item[2017] Trek Excellence Scholarship for Continuing Students\\
          University of British Columba
    \item[2016] Best Demo\\
          UBC HCI Designing for People Year-end Event
    \item[2016] Undergraduate Student Research Award\\
          Natural Sciences and Engineering Research Council of Canada (NSERC)
    \item[2016] Trek Excellence Scholarship for Continuing Students\\
          University of British Columba
    \item[2016] J. Fred Muir Memorial Scholarship\\
          University of British Columba
\end{description}

\subsubsection{Service} \hfill

11x Outstanding Reviewer Awards
\hfill	ICML (2021, 2022, 2024, 2025)

\vspace{-1ex}
\hfill NeurIPS (2020, 2022, 2023, 2025)

\vspace{-1ex}
\hfill ICLR (2022), AISTATS (2022, 2025)

TMLR Expert Reviewer \hfill Since 2024.

\vspace{2ex}

\subsection{Invited Talks and visits}
\newcommand{\where}[1]{{#1}}
\begin{description}[leftmargin=\mleftsize,style=nextline,font=\normalfont\space]
    \item[2026] Global Optimization of Deep Relu Networks Using Convex Reformulations. \hfill Endinburgh, UK \\
          SIAM Conference on Optimization (OP26)

    \item[2026] Convex Analysis of Non-Convex Neural Networks.  \hfill Oxford, UK\\
          Mathematical Institute, University of Oxford

    \item[2026] Level Set Teleportation: An Optimization Perspective.  \hfill Lausanne, Switzerland\\
          MLO Lab, EPFL

    \item[2026] Convex Analysis of Non-Convex Neural Networks. \hfill Paris, France\\
          Sierra Team, Inria

    \item[2025] Optimal Sets and Solution Paths of ReLU Networks. \hfill Berkeley, USA\\
          SIAM NCC25

    \item[2025] Level Set Teleportation: An Optimization Perspective. \hfill Montreal, Canada\\
          Montreal MLOpt Seminar, Mila

    \item[2024] Optimal Sets and Solution Paths of ReLU Networks. \hfill Lausanne, Switzerland\\
          MLO lab, EPFL

    \item[2024] Level Set Teleportation: An Optimization Perspective. \hfill Tübingen, Germany\\
          MPI Intelligent Systems and ELLIS Institute

    \item[2024] Optimal Sets and Solution Paths of ReLU Networks. \hfill Virtual\\
          Math Machine Learning Seminar MPI MIS + UCLA

    \item[2023] SGD Under Interpolation: Convergence, Line-search, and Acceleration. \hfill Seattle, USA\\
          SIAM Conference on Optimization (OP23)

    \item[2022] Fast Convex Optimization for Two-Layer ReLU Networks. \hfill Vancouver, Canada\\
          UBC Institute for Applied Mathematics (IAM)

    \item[2020] Painless Stochastic Gradient: Interpolation, Line-Search, and Convergence Rates. \hfill Virtual\\
          Machine Learning Summer School (MLSS) 202University of British Columbia\end{description}

\vspace{2ex}

\subsection{Service}

\begin{description}[leftmargin=\mleftsize,style=nextline,font=\normalfont\space]
    \item[2026] \textbf{Organizer}, Minisymposium on Global Optimization for Neural Networks \\
          SIAM Conference on Optimization (OP26)
    \item[2024] \textbf{Organizer}, ``Optimization Lunch'' Talk Series\\
          Department of Electrical Engineering, Stanford University
    \item[2019] \textbf{Organizer}, Machine Learning Reading Group (MLRG)\\
          Department of Computer Science, University of British Columbia
\end{description}

\vspace{2ex}

\subsection{Volunteering}

\begin{description}[leftmargin=\mleftsize,style=nextline,font=\normalfont\space]
    \item[2025] Stanford SERIO Mentorship Program
    \item[2025] Stanford CS Undergraduate Mentorship Program
    \item[2025] Stanford Science Penpals Program
    \item[2024] Stanford SERIO Mentorship Program
    \item[2024] Stanford CS Undergraduate Mentorship Program
    \item[2024] Stanford Student Application Support Program
\end{description}

\vspace{2ex}

\subsection{Reviewing}
\begin{description}[leftmargin=10em,style=nextline,font=\normalfont\space]
    % \begin{description}[leftmargin=10em]
    \item[NeurIPS] Neural Information Processing Systems (NeurIPS)\hfill 2019--2025
    \item[ICML] International Conference on Machine Learning (ICML) \hfill 2021--2025
    \item[ICLR] International Conference on Learning Representations \hfill 2022
    \item[AISTATS] International Conference on Artificial Intelligence and Statistics \hfill 2022--2025
    \item[TMLR] Transactions on Machine Learning Research \hfill 2023--2025
    \item[JMLR] Journal of Machine Learning Research \hfill 2021, 2022
    \item[Trans. Inf. Theory] IEEE Transactions on Information Theory \hfill 2024
    \item[SIMODS] SIAM Journal on Mathematics of Data Science \hfill 2023
    \item[OptML workshop] NeurIPS workshop: Optimization for Machine Learning \hfill 2023, 2024
\end{description}

\vspace{2ex}

\subsection{Teaching Assistantships}

\begin{itemize}[leftmargin=0pt,label={}]%,parsep=3pt,itemsep=0pt,topsep=0pt,partopsep=0pt,]
    \item EE 364B: {\it Convex Optimization II},
          Stanford University
          \hfill \mute{Spring } 2022, 2023, 2024, 2025

    \item \emph{Workshop on Bayesian Machine Learning}, Data Science Summer School (D3)
          % (Mark Schmidt)
          \hfill  2018

    \item CPSC 340: {\it Machine Learning},
          University of British Columbia 
          \hfill \mute{Fall }2017

    \item CPSC 210: {\it Software Construction},
          University of British Columbia
          \hfill \mute{Fall }2015

    \item CPSC 210: {\it Computation, Programs, and Programming},
          University of British Columbia
          \hfill \mute{Fall }2014
\end{itemize}

\vspace{1ex}

\subsection{Open Source Software}
\begin{description}[leftmargin=6.5em,style=nextline,font=\normalfont\space]
    \item[SCNN] 
          A library for global optimization of two-layer ReLU networks using
          convex reformulations.
          \\
          \texttt{\url{https://github.com/pilancilab/scnn}}
\end{description}
